\documentclass[10pt,conference]{IEEEtran}

\usepackage[T1]{fontenc}
\usepackage[utf8]{inputenc}
\usepackage{amsmath,amssymb,amsfonts}
\usepackage{graphicx}
\usepackage{booktabs}
\usepackage{array}
\usepackage{multirow}
\usepackage{xcolor}
\usepackage{url}
\usepackage{hyperref}
\usepackage{cite}
\usepackage{microtype}
\usepackage{placeins}

\hypersetup{colorlinks=true,linkcolor=blue,citecolor=blue,urlcolor=blue}
\newcommand{\safeincludegraphics}[2][]{%
  \IfFileExists{#2}{%
    \includegraphics[#1]{#2}%
  }{%
    \fbox{\parbox{0.95\columnwidth}{\centering Missing figure:\\\texttt{#2}}}%
  }%
}

\title{Game Playing with Q-Learning:\\
Tabular Reinforcement Learning for Simplified Snake}

\author{
  \IEEEauthorblockN{Group 999 --- Track A: Game Playing with Q-Learning}
  \IEEEauthorblockA{\textit{Anyi Wang, Junyi He, Yuqi Zhu, Yuqi Ren} \\
  Reinforcement Learning Course Final Project, 2026}
}

\begin{document}
\maketitle

%% -------------------------------------------------------------
\begin{abstract}
We study tabular Q-learning for a simplified Snake game on a $6{\times}6$ grid. The state space is efficiently compressed into 11 binary features encoding local danger, heading, and food direction, reducing the problem to 2,048 discrete states. Our primary $\varepsilon$-greedy Q-learning agent, trained over 10{,}000 episodes, achieves an average greedy-evaluation score of 13.73, representing a 73$\times$ improvement over the random baseline (0.19). We further evaluate the impact of exploration schedules, TD-control variants, and reward-shaping methods.  Our findings demonstrate that: (i) exploration scheduling exerts a more significant influence on performance than specific algorithm choice; (ii) potential-based shaping provides marginal benefits while simple shaping can lead to sub-optimal policies; and (iii) the compact 11-bit encoding imposes a structural performance ceiling for longer snake configurations. 
\end{abstract}

%% -------------------------------------------------------------
\section{Introduction}
\label{sec:intro}

Tabular Q-learning~\cite{watkins1992q} remains the canonical entry
point to reinforcement learning: every value is stored explicitly,
making the learned knowledge fully auditable.
Yet even in this transparent regime, subtle design choices---how fast
to decay exploration, whether to add reward shaping, which TD target
to use---produce substantial performance differences that deserve
rigorous empirical analysis.

The Snake game encapsulates four challenges that are representative
of broader game-playing RL:
\emph{sparse rewards} (food events are infrequent),
\emph{dynamic self-avoidance} (the growing body creates
time-varying obstacles),
\emph{exploration-exploitation tension}, and
\emph{partial observability} under compact state encodings.
On a $6{\times}6$ grid with an 11-bit state encoding
($|\mathcal{S}|=2{,}048$), tabular methods are not just
tractable but also fully interpretable: individual Q-values can be
decoded back to game semantics, and visit counts expose exactly
which parts of the state space the agent has explored.

This report makes four contributions.
We design the MDP and its compact state encoding
(Section~\ref{sec:formulation}).
We provide algorithmic descriptions of four TD-control methods and
three exploration strategies, grounded in theory
(Section~\ref{sec:methods}).
We perform in-depth quantitative analysis of training dynamics,
Q-value semantics, convergence behaviour, and all ablation variants,
drawing on multiple independent experimental datasets
(Section~\ref{sec:experiments}).
Finally, we discuss open research challenges
(Section~\ref{sec:challenges}).

%% -------------------------------------------------------------
\section{Problem Formulation}
\label{sec:formulation}

\subsection{MDP Specification}

We model Snake as a discrete-time MDP
$\mathcal{M} = (\mathcal{S}, \mathcal{A}, P, R, \gamma, T)$:
\begin{itemize}
  \item $\mathcal{S}$: 11-bit discrete state space
        ($|\mathcal{S}|=2{,}048$);
  \item $\mathcal{A}=\{0,1,2\}$: \textsc{straight},
        \textsc{turn-left}, \textsc{turn-right};
  \item $P(s'|s,a)$: deterministic body dynamics, stochastic food
        respawn (uniform over unoccupied cells);
  \item $R$: immediate scalar reward (see below);
  \item $\gamma=0.9$; $T=1{,}000$.
\end{itemize}

\subsection{State Representation}

The full Snake board contains the head position, ordered body layout,
food location, and current direction; treating that raw grid as the
state would lead to a combinatorial explosion that is impractical for a
small tabular method.
The 11-dimensional binary feature vector is:
\begin{equation}
  s = [\underbrace{d_s, d_l, d_r}_{\text{danger}},\;
       \underbrace{\delta_U, \delta_D, \delta_L, \delta_R}_{\text{heading (one-hot)}},\;
       \underbrace{f_U, f_D, f_L, f_R}_{\text{food direction}}],
\end{equation}
where $d_s,d_l,d_r$ indicate one-step collision danger and
$f_U,f_D,f_L,f_R$ encode the sign of the food offset relative to
the head (at most one bit per axis).
The state integer is $s=\sum_{i=0}^{10}b_i\cdot 2^{10-i}$, giving
$|\mathcal{S}|=2{,}048$, but only 256 states (12.5\%) are reachable
under valid game dynamics on this grid size.

\textbf{Partial observability limitation.}
The encoding captures only one-step look-ahead danger; body segments
beyond the immediate neighbourhood are invisible.
This creates a fundamental performance ceiling discussed in
Section~\ref{sec:ceiling}.

\subsection{Reward Function}

\begin{equation}
  R(s,a,s') =
  \begin{cases}
    +10 & \text{food eaten,}\\
    -10 & \text{collision (wall or body),}\\
    -0.1 & \text{otherwise.}
  \end{cases}
\end{equation}

The 100:1 food/step ratio ensures the dominant signal is food
collection; the step penalty prevents looping.

%% -------------------------------------------------------------
\section{Key RL Methods}
\label{sec:methods}

\subsection{Q-Learning (Primary)}

Off-policy TD control~\cite{watkins1992q}:
\begin{equation}
  Q(s,a)\leftarrow Q(s,a)+\alpha\bigl[
    r + \gamma\max_{a'}Q(s',a')\cdot\mathbf{1}[\neg\text{done}]
    -Q(s,a)\bigr].
  \label{eq:ql}
\end{equation}
Here $\alpha$ is the learning rate, $\gamma$ is the discount factor,
and the behaviour policy uses $\varepsilon$-greedy exploration:
with probability $\varepsilon$ the agent chooses a random action, and
with probability $1-\varepsilon$ it chooses the current greedy action.
All Q-values initialise to 0; ties in $\arg\max$ are broken
uniformly at random.

\subsection{SARSA and Expected SARSA}

SARSA (on-policy)~\cite{rummery1994line} bootstraps from the
actual next action:
$Q(s,a)\leftarrow Q(s,a)+\alpha[r+\gamma Q(s',a')-Q(s,a)]$.

Expected SARSA~\cite{sutton2018reinforcement} replaces the sampled
action with the policy expectation:
$\mathbb{E}_{a''\sim\pi}[Q(s',a'')]$,
reducing variance while remaining on-policy.

\subsection{Double Q-Learning}

Double Q-learning~\cite{hasselt2010double} maintains two tables
$Q_1,Q_2$, randomly updating one at each step with the other as
value evaluator:
\begin{equation}
  Q_1(s,a)\leftarrow Q_1(s,a)+\alpha\bigl[
    r+\gamma Q_2\!\bigl(s',\arg\max_{a'}Q_1(s',a')\bigr)-Q_1(s,a)\bigr].
\end{equation}
Action selection uses $Q_1+Q_2$.

\subsection{Exploration Strategies}

\textbf{$\varepsilon$-Greedy (primary).}
Random action with probability $\varepsilon$, greedy otherwise.
Schedule: $\varepsilon_0=1.0$, $\varepsilon\leftarrow
\max(\varepsilon\cdot d,0.01)$ per episode.

\textbf{Softmax.}
$\pi(a|s)\propto\exp(Q(s,a)/\tau)$, with temperature $\tau$
decaying analogously.

\textbf{UCB.}
$a^*=\arg\max_a\bigl[Q(s,a)+c\sqrt{\ln(N(s)+1)/(N(s,a)+10^{-8})}\bigr]$,
$c=1$.

\subsection{Reward Shaping}

Following Ng et al.~\cite{ng1999policy}:

\textbf{Simple shaping.}
$+w$ ($-w$) if Manhattan distance to food decreases (increases).

\textbf{Potential-based shaping.}
$F(s,s')=\gamma\Phi(s')-\Phi(s)$, $\Phi(s)=-d_{\text{food}}(s)$.
Policy-invariant by construction~\cite{ng1999policy}.

\subsection{Training Protocol}

See Table~\ref{tab:hyper} for primary hyperparameters.
Each variant is evaluated by 300 greedy ($\varepsilon=0$) episodes
post-training and compared to a 300-episode uniform-random baseline.

\begin{table}[h]
  \centering
  \caption{Primary Q-Learning Hyperparameters}
  \label{tab:hyper}
  \begin{tabular}{ll}
    \toprule
    Parameter & Value \\
    \midrule
    Grid size                       & $6\times6$ \\
    Initial snake length            & 2 \\
    Max steps / episode             & 1\,000 \\
    Training episodes               & 10\,000 \\
    Greedy evaluation episodes      & 300 \\
    $\alpha$ (learning rate)        & 0.1 \\
    $\gamma$ (discount)             & 0.9 \\
    $\varepsilon_0$                 & 1.0 \\
    $\varepsilon$ decay $d$         & 0.995 \\
    $\varepsilon_{\min}$            & 0.01 \\
    Q-init                          & 0.0 \\
    Random seed                     & 42 \\
    \bottomrule
  \end{tabular}
\end{table}

%% -------------------------------------------------------------
\section{Experimental Results and Comparisons}
\label{sec:experiments}

All reported numbers are drawn directly from experimental CSV logs
and saved \texttt{.npy} Q-table files.
We distinguish two metrics throughout: the \emph{training
moving-average (MA) score} (window 100, computed with
$\varepsilon>0$) and the \emph{greedy evaluation score} (300
episodes, $\varepsilon=0$, post-training).
These differ significantly for exploration-heavy configurations,
as we explain below.

\begin{figure}[tbp]
  \centering
  \safeincludegraphics[width=\columnwidth]{results/figures/moving_avg_score.png}
  \caption{Moving-average score trend from \texttt{results/figures/moving\_avg\_score.png}. The curve shows stable improvement of the main Q-learning run over training.}
  \label{fig:moving_avg_score}
\end{figure}

\begin{figure}[tbp]
  \centering
  \safeincludegraphics[width=\columnwidth]{results/figures/moving_avg_reward.png}
  \caption{Moving-average reward trend from \texttt{results/figures/moving\_avg\_reward.png}, consistent with the score improvement pattern.}
  \label{fig:moving_avg_reward}
\end{figure}

\begin{figure}[tbp]
  \centering
  \safeincludegraphics[width=\columnwidth]{results/figures/baseline_vs_q_learning.png}
  \caption{Greedy evaluation comparison between random baseline and trained Q-learning policy from \texttt{results/figures/baseline\_vs\_q\_learning.png}.}
  \label{fig:baseline_vs_ql}
\end{figure}

\begin{figure}[tbp]
  \centering
  \safeincludegraphics[width=\columnwidth]{results/figures/epsilon_decay_comparison.png}
  \caption{Effect of epsilon-decay settings from \texttt{results/figures/epsilon\_decay\_comparison.png}.}
  \label{fig:epsilon_decay_visual}
\end{figure}

\begin{figure}[tbp]
  \centering
  \safeincludegraphics[width=\columnwidth]{results/figures/reward_shaping_comparison.png}
  \caption{Reward-shaping comparison from \texttt{results/figures/reward\_shaping\_comparison.png}.}
  \label{fig:reward_shaping_visual}
\end{figure}

\subsection{Primary Result: Trained Agent vs.\ Random Baseline}

\begin{table}[h]
  \centering
  \caption{Trained Q-Learning vs.\ Random Baseline (300 greedy episodes)}
  \label{tab:primary}
  \begin{tabular}{lccccc}
    \toprule
    Agent & Avg.\ Score & Std & Max & Avg.\ Reward & Avg.\ Steps\\
    \midrule
    Q-Learning & \textbf{13.73} & 4.29 & \textbf{27} & \textbf{121.67} & \textbf{70.0}\\
    Random     & 0.19           & 0.44 &  3  & $-$8.89  & 8.7\\
    Ratio      & \textbf{72.8$\times$} & --- & \textbf{9$\times$} & ---
               & \textbf{8.0$\times$}\\
    \bottomrule
  \end{tabular}
\end{table}

The 72.8$\times$ score improvement (Table~\ref{tab:primary}) is the
headline result. Several secondary observations are equally
informative.

\textbf{Max score of 27 on a $6{\times}6$ grid.}
A score of 27 means the snake occupies $2+27=29$ of 36 cells
(80.6\%).
At this density the board is nearly saturated, and valid routes to
food exist only through narrow corridors.
Reaching this score requires the agent to navigate without any
information about deep body segments---a non-trivial feat under the
11-bit encoding.

\textbf{Survival steps: 70.0 vs 8.7.}
The 8$\times$ survival improvement implies the agent completes
$\approx70/5.1\approx13.7$ food segments per episode on average
(consistent with score 13.73).
The random agent dies within 9 steps almost exclusively through
collision, confirming zero collision-avoidance learning.

\textbf{Score variance.}
Standard deviation 4.29 (coefficient of variation $CV=0.31$) reflects
the structural partial-observability of the encoding: even the mature
policy occasionally walks into multi-step body traps it cannot
foresee (Section~\ref{sec:ceiling}).

\subsection{Four-Phase Training Dynamics}
\label{sec:phases}

Analysis of the 10\,000-episode \texttt{basic\_learning.csv} log
reveals four qualitatively distinct learning phases
(Table~\ref{tab:phases}).

\begin{table}[h]
  \centering
  \caption{Four Training Phases (100-episode MA window)}
  \label{tab:phases}
  \begin{tabular}{cccccc}
    \toprule
    Phase & Episodes & Avg Score & Std & Zero\% & $\bar{\varepsilon}$ \\
    \midrule
    I   & 1--500    & 2.87  & 2.77 & 20.8\% & 0.366 \\
    II  & 501--1000 & 9.94  & 4.44 &  0.2\% & 0.030 \\
    III & 1001--3000 & 11.65 & 4.63 & 0.1\% & 0.010 \\
    IV  & 3001--10000 & 12.54 & 4.57 & 0.2\% & 0.010 \\
    \bottomrule
  \end{tabular}
\end{table}

\textbf{Phase~I (eps.\ 1--500): Random exploration and Q-table seeding.}
$\varepsilon$ decays from 1.0 to 0.082 by episode 500.
The agent is almost entirely random; 20.8\% of episodes score zero
(immediate death with no food collected).
Yet this phase is crucial: it seeds the Q-table from scratch.
The MA score crosses 5.0 at episode 463 when $\varepsilon=0.098$,
precisely at the boundary where exploitative actions begin to
dominate---confirming that early Q-table data first produces
a usable signal exactly when the agent starts following its values.

\textbf{Phase~II (eps.\ 501--1000): Rapid policy crystallisation.}
This is the steepest learning interval.
$\varepsilon$ falls from 0.082 to 0.010 (the minimum);
the MA score surges from $\approx$6.0 to $\approx$11.2,
crossing 8.0 at episode 643 ($\varepsilon=0.040$) and
10.0 at episode 716 ($\varepsilon=0.028$).
Score standard deviation peaks at 4.44: the policy is in
transition, with some episodes already near-optimal while others
still reflect partially converged Q-values.
The zero-episode fraction drops from 20.8\% to 0.2\%, confirming
that collision avoidance is learned almost completely in this phase.

\textbf{Phase~III (eps.\ 1001--3000): Convergence and fine-tuning.}
With $\varepsilon$ locked at 0.010, the policy is 99\% greedy.
MA score stabilises around 11.65--12.31.
For the main Q-learning run, the number of visited states saturates at
exactly 256 by episode $\approx$2\,000 and does not grow thereafter.

\textbf{Phase~IV (eps.\ 3001--10000): Plateau.}
MA score oscillates between 12.3 and 12.6 without a sustained
upward trend.
The maximum episode score achieved during training reaches
\textbf{35} (vs 27 in greedy evaluation), occurring in late training
due to lucky body configurations that happen not to create traps.
All episodes in the last 1\,000 terminate by collision; none
reaches the 1\,000-step cap, confirming that the snake never
fills the board under the current policy.

\subsection{Q-Value Convergence and Semantic Interpretation}
\label{sec:qvalue}

After 611\,892 total updates:
$Q_{\max}=23.88$, $Q_{\min}=-10.0$, $\bar{Q}=+0.271$,
$|Q|_{\text{mean}}=0.936$, $\sigma_Q=3.185$.

\textbf{Q-value distribution structure.}
Of the 6\,144 Q-table entries, 87.9\% are exactly zero
(never updated), 6.0\% are positive ($>0.01$), and
6.1\% are negative ($<-0.01$).
The zero entries correspond to the 87.5\% of states
never visited.
The symmetry between positive (6.0\%) and negative (6.1\%)
entries reflects the balance between food-seeking (positive
expected return) and collision-risk (negative expected return)
in the reachable state space.

\textbf{Semantic decoding of highest Q-values.}
Table~\ref{tab:topq} decodes the five highest Q-table entries
back into game semantics.

\begin{table}[h]
  \centering
  \caption{Top-5 Q-Table Entries: State Semantics}
  \label{tab:topq}
  \begin{tabular}{clcl}
    \toprule
    Q-value & Action & Danger bits & Heading / Food \\
    \midrule
    23.88 & Straight & right only    & Facing LEFT, food LEFT \\
    21.75 & Straight & none          & Facing UP, food UP \\
    21.72 & Turn-left & none         & Facing RIGHT, food UP \\
    21.59 & Turn-left & straight     & Facing RIGHT, food UP \\
    21.57 & Turn-left & straight+right & Facing UP, food LEFT \\
    \bottomrule
  \end{tabular}
\end{table}

These entries reveal clear policy logic.
The highest Q-value (23.88) occurs when the snake faces LEFT with
food directly ahead (LEFT) and a danger only on the right:
the correct action is to go straight toward food, which the
Q-table correctly identifies.
The fourth entry ($Q=21.59$, state 1048) is particularly
informative: the snake faces RIGHT, food is UP (requiring a
left turn), and going straight is blocked---the agent correctly
identifies turn-left as the escape action.
The fifth entry ($Q=21.57$, state 1410) shows the agent choosing
turn-left when both straight and right are dangerous and food is
to the left---again the unique feasible move.
These semantically coherent high-Q states confirm that the learned
policy is not a statistical artefact but encodes genuine
navigation knowledge.

\textbf{Dangerous states: when all actions fail.}
Inspecting entries with $Q_{\max}<-5$ and visit count $>0$
reveals 32 states in which every reachable action leads
near-certainly to death.
These states all share $d_s=d_l=d_r=1$ (all three one-step
actions are immediately dangerous), representing the snake
trapped in a corner or hemmed in by its own body.
The lowest entry is $Q(s=1814,\text{all})=-10.0$ with 335 visits:
the agent has learned that this configuration is a dead end, but
cannot avoid it due to the look-ahead limitation.

\textbf{Residual TD error.}
The mean absolute TD error stabilises at 2.57 in the last 1\,000
episodes, not at zero, even though Q-values appear converged.
This irreducible residual arises from food respawning stochasticity:
the same $(s,a,r,s')$ tuple occurs with different next-state
Q-values depending on where food lands.
Mathematically, the Bellman operator $\mathcal{T}Q(s,a)=
\mathbb{E}[r+\gamma\max_{a'}Q(s',a')]$ has a variance due to
$s'$ randomness that prevents the empirical TD error from reaching
zero even at true convergence.
Notably, TD error is slightly \emph{higher} in high-score episodes
($\geq21$: 2.59 vs $0$--$5$: 2.40) because longer trajectories
generate more diverse state transitions within a single episode.

\subsection{Structural Performance Ceiling}
\label{sec:ceiling}

The training MA score never exceeds 13 during 10\,000 episodes
(Figure not shown due to PDF constraints; confirmed by data).
We identify this as a \textbf{structural limit} imposed by the
11-bit encoding, supported by three independent lines of evidence:

\begin{enumerate}
  \item \textbf{State saturation.}
        Visited states plateau in a very narrow range across the full
        ablation suite: most variants end at 256 visited states, while
        Softmax reaches 257 and slow-decay reaches 259.
        No further training discovers new states; improvement must come
        from refining already-visited Q-values,
        which yields diminishing returns after Phase~II.

  \item \textbf{Score distribution compression.}
        In the last 2\,000 training episodes of the baseline variant,
        44.3\% of scores fall in $[10,14]$, 25.4\% in $[15,19]$, and
        only 0.4\% exceed 24 (Table~\ref{tab:scoredist}).
        The distribution is truncated at the upper end not by episode
        length (the 1\,000-step cap is never reached) but by the agent
        encountering body traps invisible to the 11-bit state.

  \item \textbf{One-step danger horizon.}
        Decoding state semantics confirms that all Q-table knowledge
        is limited to one-step danger ($d_s,d_l,d_r$).
        When the snake reaches score $\geq15$, non-immediate body
        segments create two- and three-step traps that all map to the
        same one-step-safe state encoding,
        leading to seemingly random deaths from the agent's perspective.
\end{enumerate}

\FloatBarrier
\begin{table*}[!t]  % use table* and place at top of page
  \centering
  \caption{Score Frequency Distribution (last 2\,000 training episodes)}
  \label{tab:scoredist}
  \begin{tabular}{lccccccc}
    \toprule
    Variant & [0--4] & [5--9] & [10--14] & [15--19] & [20--24] & [25--29] & [$\geq$30]\\
    \midrule
    QL-$\varepsilon$-greedy & 4.7\% & 19.6\% & 44.3\% & 25.4\% &  5.5\% & 0.4\% & 0.1\%\\
    QL-slow-decay     & 4.7\% & 18.4\% & 42.5\% & 27.2\% &  6.3\% & 0.8\% & 0.1\%\\
    QL-softmax        & 1.8\% & 20.1\% & 42.6\% & 30.4\% &  4.7\% & 0.3\% & 0.0\%\\
    QL-UCB            & 1.1\% & 20.0\% & 42.5\% & 31.4\% &  4.9\% & 0.1\% & 0.0\%\\
    Double-Q          & 3.9\% & 17.6\% & 42.2\% & 30.2\% &  5.6\% & 0.5\% & 0.0\%\\
    QL-potential      & 4.3\% & 17.8\% & 42.0\% & 28.7\% &  6.7\% & 0.5\% & 0.1\%\\
    QL-simple shaping & 4.5\% & 23.4\% & 41.3\% & 25.8\% &  4.5\% & 0.4\% & 0.0\%\\
    \bottomrule
  \end{tabular}
\end{table*}

\subsection{Full Ablation Study}
\label{sec:ablation}

\begin{table*}[!t]
  \centering
  \caption{Full Ablation (10\,000 training episodes,
           300 greedy evaluation episodes)}
  \label{tab:ablation}
  \begin{tabular}{llllcccccc}
    \toprule
    Variant & Algorithm & Exploration & Shaping
            & Greedy & Std & Max & $|Q|_{\text{mean}}$
            & Total visits & Pairs updated \\
    \midrule
    Slow $\varepsilon$-decay & Q-learning & $\varepsilon$-greedy ($d{=}0.999$) & none
        & \textbf{14.79} & 4.28 & \textbf{30} & 1.056 & 514\,951 & 745 \\
    Double Q-learning & Double Q   & $\varepsilon$-greedy & none
        & 13.98 & 4.31 & 25 & 1.638 & 632\,599 & 744 \\
    Potential shaping & Q-learning & $\varepsilon$-greedy & potential
        & 13.80 & \textbf{3.97} & 27 & 0.964 & 619\,410 & 743 \\
    \textbf{Q-Learning (main)} & Q-learning & $\varepsilon$-greedy & none
        & 13.73 & 4.29 & 27 & 0.936 & 611\,892 & 742 \\
    Expected SARSA & Exp.\ SARSA & $\varepsilon$-greedy & none
        & 13.45 & 4.34 & \textbf{30} & 0.911 & 597\,561 & 750 \\
    UCB ($c{=}1$) & Q-learning & UCB & none
        & 13.42 & 4.06 & 26 & 0.625 & 684\,016 & 768 \\
    Softmax & Q-learning & Softmax & none
        & 13.37 & 4.12 & 29 & 0.635 & 638\,476 & 694 \\
    SARSA & SARSA & $\varepsilon$-greedy & none
        & 13.33 & 4.01 & 24 & 0.931 & 596\,060 & 747 \\
    Simple shaping & Q-learning & $\varepsilon$-greedy & simple
        & 12.99 & 4.32 & 27 & 0.957 & 576\,999 & 744 \\
    Low $\varepsilon_0$ ($=0.5$) & Q-learning & $\varepsilon$-greedy & none
        & 12.79 & 4.45 & 28 & 0.884 & 590\,495 & 741 \\
    \midrule
    Random baseline & --- & uniform & ---
        & 0.19 & 0.44 & 3 & --- & --- & --- \\
    \bottomrule
  \end{tabular}
\end{table*}

\subsubsection{Exploration Schedule: The Dominant Factor}

The \texttt{epsilon\_decay\_comparison.csv} dataset provides a
controlled three-way comparison with the same algorithm but
different exploration schedules
($\varepsilon_0,d$):
slow ($1.0,\,0.999$), medium ($1.0,\,0.995$), and
low-exploration ($0.3,\,0.995$).
Table~\ref{tab:epsdecay} shows training-MA milestones.

\begin{table}[h]
  \centering
  \caption{Training MA Score at Key Episodes: Epsilon Strategies}
  \label{tab:epsdecay}
  \begin{tabular}{lcccccc}
    \toprule
    Strategy & 100 & 300 & 500 & 1000 & 5000 & 10000 \\
    \midrule
    Slow ($d{=}0.999$)        & 0.28 & 0.66 & 1.15 & 2.32 & 12.75 & 12.16 \\
    Medium ($d{=}0.995$)      & 0.49 & 2.62 & 6.02 & 11.20 & 12.58 & 12.31 \\
    Low ($\varepsilon_0{=}0.3$) & 1.58 & 6.30 & 9.30 & 11.25 & 12.75 & 11.84 \\
    \bottomrule
  \end{tabular}
\end{table}

The table reveals a striking \emph{speed-versus-quality trade-off}.
Low-exploration learns fastest in episodes 1--500 (MA=9.30 vs 6.02
for medium), but plateaus earliest and finishes lowest (11.84).
Slow decay is slowest in training (MA=1.15 at ep500) but ultimately
achieves a comparable training MA (12.16) to medium (12.31) at
episode 10\,000---and the highest \emph{greedy evaluation score}
(14.79) in the ablation.

The key insight is the \textbf{training-vs-evaluation discrepancy}.
Slow decay ($d=0.999$) does not reach $\varepsilon=0.01$ until
episode 4\,508; the agent is still exploring ($\varepsilon\approx0.03$)
at episode 5\,000.
Its training MA score of 12.16 includes the drag of exploratory
actions.
When evaluated greedily ($\varepsilon=0$) post-training, it achieves
14.79---because it received far more diverse TD updates before
committing to exploitation.
Medium decay reaches $\varepsilon_{\min}$ at episode 900, after
which no further Q-table diversification occurs.

Quantitatively: slow decay records 514\,951 state visits over
10\,000 episodes vs 611\,892 for the medium schedule, yet it still
achieves a 7.7\% higher greedy score, suggesting that
\emph{update quality matters more than update quantity} in this
tabular setting.

\subsubsection{Convergence Speed: Softmax and UCB Are Faster}

A counterintuitive finding from Table~\ref{tab:conv} is that
Softmax and UCB converge to high performance \emph{much faster}
than $\varepsilon$-greedy in early training.

\begin{table}[h]
  \centering
  \caption{Episode at which Training MA Score First Crosses Threshold}
  \label{tab:conv}
  \begin{tabular}{lcccc}
    \toprule
    Variant & MA$\geq$5 & MA$\geq$8 & MA$\geq$10 & MA$\geq$12 \\
    \midrule
    QL-UCB      & \textbf{171} & \textbf{250} & \textbf{305} & \textbf{360} \\
    QL-Softmax  & 194 & 250 & 303 & 390 \\
    QL-Pot.\ shaping & 436 & 582 & 678 & 842 \\
    QL-$\varepsilon$-greedy & 463 & 643 & 716 & 914 \\
    SARSA       & 421 & 590 & 740 & 930 \\
    Exp.\ SARSA & 421 & 614 & 739 & 981 \\
    Double-Q    & 445 & 560 & 705 & 913 \\
    QL-slow     & 1666 & 2556 & 3067 & 3683 \\
    QL-low-$\varepsilon_0$ & \textit{372} & 538 & 735 & 7833 \\
    \bottomrule
  \end{tabular}
\end{table}

UCB crosses MA=10 at episode 305, Softmax at 303---a
\textbf{2.4$\times$ speedup} over $\varepsilon$-greedy (episode 716).
At episode 300, UCB achieves MA=9.72 vs $\varepsilon$-greedy's 2.62.
These early advantages arise from the richer exploration signals:

Softmax assigns graded probabilities proportional to Q-values from
the first episode, immediately concentrating experience on
higher-valued actions while retaining some probability for all
actions.
Even when Q-values are near-zero (early training), Softmax exploration
is equivalent to uniform with a slight push toward actions with any
positive value signal.
UCB's optimism bonus systematically tries each action in every
newly visited state, guaranteeing at least one update per
action per state before settling into value-based selection.
This ensures faster Q-table coverage in early episodes.

However, this early advantage does not translate to better
long-run performance.
By episode 10\,000, $\varepsilon$-greedy (greedy eval: 13.73)
outperforms Softmax (13.37) and UCB (13.42) in greedy evaluation.
This reversal occurs because:
(i) at $\varepsilon_{\min}=0.01$, all three strategies are essentially
equivalent in exploitation;
(ii) UCB's count-based bonus continues to ``waste'' some updates on
under-visited pairs even in late training, reducing the quality of
updates in high-value regions;
(iii) Softmax visits only 694 state-action pairs (vs 742 for
$\varepsilon$-greedy), focusing experience on currently promising
states at the cost of less overall coverage.

\subsubsection{On-Policy vs.\ Off-Policy: Algorithm Comparison}

Q-learning (13.73) $>$ Expected SARSA (13.45) $>$ SARSA (13.33).
The off-policy advantage over SARSA is 0.40 points (3.0\%).

Under $\varepsilon$-greedy with $|\mathcal{A}|=3$ and
$\varepsilon=0.01$, the on-policy bootstrap target differs from
the off-policy target by:
$\Delta = (1-\varepsilon)(0) + \varepsilon/3 \cdot
[\bar{Q}(s')-\max_{a'}Q(s',a')]$.
Since $\bar{Q}(s')\leq\max Q(s',a')$, this is always non-positive,
meaning SARSA systematically bootstraps from a slightly pessimistic
target.
On a small board where survival is paramount, this conservatism
marginally reduces the value estimates in states close to danger.

Expected SARSA falls between the two: its expectation over the
$\varepsilon$-greedy policy reduces variance relative to SARSA while
still being on-policy.
Notably, Expected SARSA and SARSA \emph{update the most}
state-action pairs (750 and 747 respectively), slightly more than
Q-learning (742).
This is because on-policy training concentrates experience in states
the current policy actually visits, which---especially early in
training---includes more diverse transitions than the greedy path.

The highest maximum scores across all variants are achieved by
Expected SARSA and Slow-decay Q-learning (both 30).
Expected SARSA's lower-variance updates appear to produce
occasionally more consistent navigation in high-density
board configurations, even though its mean score is lower.

\subsubsection{Double Q-Learning: Bias, Correlation, and Scale}

Double Q-learning achieves 13.98, outperforming standard Q-learning
(13.73) by 1.8\%.
The improvement is consistent with bias reduction theory but its
magnitude reveals additional nuance.

\textbf{Q-table correlation.}
After 10\,000 episodes, $Q_1$ and $Q_2$ have Pearson correlation
$r=0.983$.
Despite being updated independently on random 50/50 splits, the two
tables converge to nearly identical estimates.
This high correlation confirms that the tabular setting with only
6\,144 entries provides sufficient coverage for both tables to
converge toward the same $Q^*$, rather than maintaining independent
uncertainty estimates as the theory ideally envisions.

\textbf{Overestimation check.}
To quantify residual overestimation, we compare the mean
maximum-action Q-value per state.
For single Q-learning: $\frac{1}{256}\sum_s\max_a Q(s,a)=1.143$.
For Double Q-learning (scaled by $\tfrac{1}{2}$ to match single
table scale): $\frac{1}{256}\sum_s\max_a\frac{Q_1+Q_2}{2}(s,a)=1.176$.
The Double Q scaled mean-max is slightly \emph{higher} than single Q,
so this metric alone does not provide strong empirical evidence of a
large overestimation reduction in this repository.
Instead, the main defensible empirical conclusion is narrower:
Double Q-learning attains a slightly higher greedy score
(13.98 vs 13.73) while keeping lower late-stage TD error
(Table~\ref{tab:tderr}).

\textbf{Scale artefact.}
The effective Q-table ($Q_1+Q_2$) shows $|Q|_{\text{mean}}=1.638$
vs 0.936 for single Q---a 1.75$\times$ ratio---and
$Q_{\max}=45.23$ vs 23.88.
This is a pure scale artefact: action selection uses the sum of two
tables, doubling all values.
The Q-min of $-20.0$ is likewise double the collision penalty of
$-10$.
Practitioners should rescale by $\tfrac{1}{2}$ when comparing
Double Q values to single-table Q-values.

\subsubsection{Action Distribution: What the Policy Learned}

Table~\ref{tab:actions} shows the action distribution across all
state visits during training.

\begin{table}[h]
  \centering
  \caption{Action Distribution Over All Training Steps}
  \label{tab:actions}
  \begin{tabular}{lccc}
    \toprule
    Variant & Straight & Left & Right \\
    \midrule
    Q-Learning        & 45.1\% & 33.1\% & 21.9\% \\
    Slow-decay        & \textbf{51.1\%} & 22.8\% & 26.1\% \\
    Low-$\varepsilon_0$ & 40.3\% & 23.7\% & 36.0\% \\
    Softmax           & 42.0\% & 28.6\% & 29.4\% \\
    UCB               & 41.4\% & 34.1\% & 24.5\% \\
    SARSA             & 44.0\% & 30.2\% & 25.8\% \\
    Expected SARSA    & 46.5\% & 32.7\% & 20.8\% \\
    Double Q          & 39.8\% & 24.1\% & \textbf{36.1\%} \\
    Simple shaping    & 42.8\% & 30.9\% & 26.3\% \\
    Potential shaping & 45.2\% & 33.1\% & 21.7\% \\
    \bottomrule
  \end{tabular}
\end{table}

\textbf{Straight-action dominance.}
All variants favour \textsc{straight} (39.8--51.1\%),
consistent with the game structure: a snake moving toward food
in a straight line is the default efficient behaviour.

\textbf{Expected SARSA is the most conservative.}
With the highest straight rate (46.5\%) and lowest turning
frequency (left 32.7\%, right 20.8\%), Expected SARSA exhibits
the most cautious policy.
Its on-policy update target includes the expected value under
the current $\varepsilon$-greedy policy, which---when the
greedy action is straight---assigns most weight to straight
and small weight to turns.
This creates a reinforcing feedback that biases the policy further
toward straight-ahead navigation.

\textbf{Double Q has the highest turning frequency (60.2\%).}
This is surprising: Double Q should produce policies similar to
single Q-learning but with less overestimation.
The high turning frequency may reflect the fact that each individual
table ($Q_1$ or $Q_2$) receives only half the updates per step,
leading to noisier individual Q-value estimates and more variable
action selection during the exploratory phase.

\textbf{Slow-decay: most straight.}
Slow-decay has the highest straight proportion (51.1\%) with the
lowest total visit count (514\,951 steps), because with higher
$\varepsilon$ during long exploration, the agent more often takes
random turns; its \emph{greedy} policy (tested during evaluation)
is more straight-biased, as reflected in its superior greedy
evaluation score (14.79).

\subsubsection{Reward Shaping: Training Signal vs.\ Policy Quality}

We have two independent data sources for reward shaping:
(1) the \texttt{reward\_shaping\_comparison.csv} experiment,
comparing sparse vs.\ shaped (simple) reward on 10\,000 episodes;
(2) the ablation experiments with simple and potential-based shaping.

\textbf{Training MA scores.}
From \texttt{reward\_shaping\_comparison.csv}:
shaped training consistently outperforms sparse training
throughout (Table~\ref{tab:shaping}).

\begin{table}[h]
  \centering
  \caption{Sparse vs.\ Shaped: Training-MA Score Milestones}
  \label{tab:shaping}
  \begin{tabular}{lccccc}
    \toprule
    Mode    & 300 & 500 & 1000 & 5000 & 10000 \\
    \midrule
    Sparse  & 2.62 & 6.02 & 11.20 & 12.58 & 12.31 \\
    Shaped  & 3.27 & 6.96 & 12.16 & 12.67 & \textbf{13.04} \\
    \bottomrule
  \end{tabular}
\end{table}

Shaping improves the MA score at ep500 by 15.6\% (6.96 vs 6.02)
and at ep1000 by 8.6\% (12.16 vs 11.20).
The faster convergence occurs because the shaping bonus
provides reward signal on \emph{every step} (positive when
moving toward food, negative otherwise), rather than only
at the rare food-eating events.

\textbf{Greedy evaluation divergence.}
Despite its better training MA, simple shaping achieves a
\emph{lower} greedy evaluation score (12.99 vs 13.73).
The source of this contradiction is revealed by analysing the
shaping bonus in the last 500 training episodes:
simple shaping adds a mean bonus of $+2.48$ per episode,
with standard deviation 1.12.
The agent has effectively learned to optimise
``reward + shaping bonus'', not just reward.
When the shaping bonus is removed at greedy evaluation,
the agent has learned a policy biased toward moving toward food
\emph{even when detours are necessary to avoid body traps},
resulting in worse actual survival.

This is a direct proxy-mismatch effect: the agent improves the
shaped training objective without improving the final greedy policy.

\textbf{Potential-based shaping escapes this problem.}
Potential shaping adds a mean bonus of $+1.97$ per episode
(std 0.99) in late training, slightly smaller than simple shaping.
Its greedy evaluation score (13.80) exceeds the no-shaping baseline
(13.73), confirming the theoretical guarantee~\cite{ng1999policy}:
the shaped reward preserves the optimal policy invariance.

The potential shaping variant also achieves the \emph{lowest}
score standard deviation in greedy evaluation (3.97 vs 4.29
for no shaping), indicating more consistent performance.
We attribute this to the smoother learning landscape: the potential
function $\Phi(s)=-d_{\text{food}}(s)$ provides a coherent
``gradient'' toward food across the entire state space,
helping the agent build a more stable value function.

\subsubsection{Convergence of TD Error: A Stability Diagnostic}

Table~\ref{tab:tderr} shows mean absolute TD error in the last
1\,000 training episodes.

\begin{table}[h]
  \centering
  \caption{TD Error in Last 1\,000 Training Episodes}
  \label{tab:tderr}
  \begin{tabular}{lcc}
    \toprule
    Variant & Mean TD Error & Std \\
    \midrule
    Double Q-learning          & \textbf{2.459} & 0.344 \\
    UCB                        & 2.530 & 0.324 \\
    Potential shaping          & 2.533 & 0.388 \\
    Expected SARSA             & 2.568 & 0.393 \\
    Q-Learning (main)          & 2.570 & 0.380 \\
    Softmax                    & 2.619 & 0.324 \\
    Slow decay                 & 2.619 & 0.398 \\
    SARSA                      & \textbf{2.713} & \textbf{0.468} \\
    \bottomrule
  \end{tabular}
\end{table}

SARSA has the highest TD error (2.713) and standard deviation (0.468).
On-policy updates with $\varepsilon=0.01$ occasionally sample
exploratory actions with poor Q-values, causing high-variance
targets.
SARSA's max TD error reaches 7.51---far above the 5.43 maximum
for Q-learning---confirming that exploratory action noise propagates
into the on-policy bootstrap target.

Double Q-learning achieves the lowest TD error (2.459),
consistent with its bias-reduction mechanism:
decorrelating action selection and value evaluation reduces
the variance of the max-operator, producing more stable targets.

\subsubsection{Summary: Effect-Size Ranking}

Ranking all ablation factors by their absolute impact on greedy
evaluation score (relative to the main Q-learning baseline at 13.73):

\begin{enumerate}
  \item \textbf{Slow $\varepsilon$-decay} ($d=0.999$): $+1.06$ ($+7.7\%$)
  \item \textbf{Double Q-learning}: $+0.25$ ($+1.8\%$)
  \item \textbf{Potential shaping}: $+0.07$ ($+0.5\%$)
  \item \textbf{Algorithm: Q vs.\ SARSA}: $-0.40$ ($-2.9\%$, SARSA lower)
  \item \textbf{Simple shaping}: $-0.74$ ($-5.4\%$)
  \item \textbf{Low initial $\varepsilon$} ($\varepsilon_0=0.5$): $-0.94$ ($-6.8\%$)
\end{enumerate}

Exploration scheduling choices (decay rate and initial value)
dominate all algorithm choices.
The gap between the best ($+7.7\%$) and worst ($-6.8\%$) exploration
configurations (14.8 points combined) far exceeds the gap between
on-policy and off-policy algorithms (0.40 points).

\subsection{6{\times}6 vs.\ 8{\times}8 Grid-Size Comparison}

To verify whether the main conclusion transfers to a larger board, we
compare Q-learning under the same core setup on 6{\times}6 and 8{\times}8
using \texttt{results/grid\_size\_comparison/grid\_size\_summary.csv}.

\begin{table}[h]
  \centering
  \caption{Grid-size comparison (Q-learning, 300 greedy eval episodes)}
  \label{tab:grid_size_comparison}
  \begin{tabular}{lcccc}
    \toprule
    Grid & Avg.\ Score & Normalized Score & Max Score & Avg.\ Steps \\
    \midrule
    6{\times}6 & 13.73 & 0.404 & 27 & 69.99 \\
    8{\times}8 & 16.88 & 0.272 & 36 & 111.03 \\
    \bottomrule
  \end{tabular}
\end{table}

The trained policy outperforms random baseline on both board sizes.
On 8{\times}8, raw score and survival steps increase because the board
provides more free space, while capacity-normalized performance
decreases (0.272 vs.\ 0.404), indicating that long-horizon planning
becomes harder on larger grids.

\subsection{Q-Learning-Only Optimization}

We added a Q-learning-only optimization study with fixed algorithm
(\texttt{q\_learning}) and fixed epsilon-greedy action selection.
Only Q-learning training choices are varied and compared:
learning rate, epsilon decay, reward shaping, optimistic
initialization, and safe action masking. The discount factor
$\gamma$ is held constant at 0.9 to isolate these optimization factors.

\begin{table}[h]
  \centering
  \caption{Q-learning-only optimization comparison (from \texttt{optimization\_summary.csv})}
  \label{tab:q_learning_optimization}
  \begin{tabular}{lcc}
    \toprule
    Variant (focus) & Avg.\ Score & Max Score \\
    \midrule
    Baseline & 13.73 & 27 \\
    \texttt{low\_alpha\_slow\_decay} (low-$\alpha$ + slow decay) & \textbf{14.33} & \textbf{30} \\
    Safe action mask only & 13.78 & 27 \\
    Potential shaping only & 13.78 & 26 \\
    Optimistic init + slow decay & 13.40 & 25 \\
    \bottomrule
  \end{tabular}
\end{table}

In Table~\ref{tab:q_learning_optimization}, the ``Variant (focus)''
labels are shortened names that map to CSV variants
\texttt{baseline}, \texttt{low\_alpha\_slow\_decay},
\texttt{safe\_action\_mask}, \texttt{potential\_shaping}, and
\texttt{optimistic\_slow\_decay}.

\begin{table}[h]
  \centering
  \caption{Factor-specific comparisons inside Q-learning-only optimization}
  \label{tab:q_learning_factor_breakdown}
  \begin{tabular}{lcc}
    \toprule
    Factor & Example variants tested & Observation \\
    \midrule
    Learning rate & $\alpha \in \{0.10,0.075,0.05,0.03\}$ & $\alpha{=}0.05$ best in this sweep \\
    Epsilon decay & decay $\in \{0.995,0.9985,0.999,0.9995\}$ & 0.999 best with tuned $\alpha$ \\
    Reward shaping & none vs potential shaping weight (0.05/0.1) & no gain over best no-shaping setup \\
    Optimistic init & init 0.0 vs 2.0 & optimistic init underperforms baseline \\
    Safe action mask & off vs on & small gain alone, no gain in best tuned setup \\
    \bottomrule
  \end{tabular}
\end{table}

Based on the optimization experiment, the most effective configuration
is \textbf{low learning rate + slow epsilon decay}
($\alpha{=}0.05$, epsilon decay $=0.999$, with $\gamma$ fixed at 0.9) without extra shaping or
optimistic initialization. This setting gives the highest greedy
evaluation average score (14.33) and ties for the best max score (30),
so we select it as the final Q-learning-only optimized configuration.

%% -------------------------------------------------------------
\section{Open Research Challenges}
\label{sec:challenges}

\subsection{Breaking the Tabular Ceiling via Richer Encodings}

Our experiments establish that all variants saturate in a very narrow
reachable-state range (256--259 states) and in a greedy-score range of
$\approx$13--15, limited by the one-step danger horizon.
Incrementally, two-step look-ahead bits ($d_{ss}$, etc.) could
extend the table to $2^{14}=16{,}384$ entries and may push the
ceiling to score 20--25.
For larger grids ($\geq10\times10$), DQN~\cite{mnih2015human}
with convolutional feature extraction becomes necessary.
The key DQN contributions---target network and replay buffer---
address the instability and sample correlation issues that arise
when the Q-function is parameterised by a neural network,
issues absent in the tabular regime studied here.

\subsection{Convergence Speed: Hybrid Exploration Schedules}

Our results show Softmax and UCB converge 2.4$\times$ faster
than $\varepsilon$-greedy (episode 303 vs 716 to MA=10),
but ultimately match or underperform in long-run greedy evaluation.
A natural follow-up is a \emph{hybrid schedule}: use UCB or Softmax
exploration in the first 500 episodes for fast Q-table seeding,
then switch to slow $\varepsilon$-decay for deep exploitation.
This would combine the best of both worlds---fast early convergence
and high final policy quality.

\subsection{Reward Shaping: Learning the Potential Function}

Our experiment confirms that the quality of the potential function
is decisive: a well-designed potential (negative Manhattan distance)
helps, while a naive simple bonus hurts.
However, the Manhattan distance potential ignores body obstacles
that block direct paths to food.
A body-aware potential function---e.g., based on flood-fill
reachability---would provide a more accurate signal for
long-snake configurations.
The challenge is computing such a potential at each step without
excessive overhead.
Alternatively, the potential function could be learned from expert
demonstrations or planning heuristics, but this would require extra
supervision not present in the current repository.

\subsection{Partial Observability and Memory}

The 32 trapped states (all actions dangerous) and the performance
ceiling at score $\approx$15 both trace to partial observability.
The POMDP framework~\cite{kaelbling1998planning} formalises this,
but belief-state planning is intractable for large state spaces.
A practical direction is to augment the 11-bit encoding with
$k$-step body history via frame stacking, or to use an LSTM-based
policy~\cite{kapturowski2018recurrent} that maintains a compressed
memory of recent body configurations.

\subsection{Non-Stationarity and the Growing-Snake Problem}

As the snake grows from length 2 to 20+, the effective dynamics change:
body cells that were irrelevant at length 2 become critical obstacles
at length 20.
The Q-table must adapt to this non-stationarity.
Curriculum learning---gradually increasing initial length during
training---could force the agent to learn body-aware navigation
before the body becomes critical.

\subsection{Exploration in Larger State Spaces}

On the $6{\times}6$ grid, all three exploration strategies cover
the same 256 reachable states.
On larger boards, the reachable state space grows combinatorially,
and random exploration becomes exponentially less efficient.
Count-based intrinsic motivation~\cite{bellemare2016unifying} and
curiosity-driven methods~\cite{pathak2017curiosity} provide principled
solutions, but their calibration to the game's extrinsic reward
timescale remains an open problem.

%% -------------------------------------------------------------
\section{Conclusion}
\label{sec:conclusion}

We have presented a comprehensive study of tabular Q-learning for
simplified Snake, grounded in direct analysis of experimental CSV
logs and saved Q-table arrays.

Our primary trained agent achieves a greedy-evaluation average score
of 13.73 (max: 27), a 73-fold improvement over the random baseline.
Through in-depth analysis, we established the following findings
beyond a simple score table:

\textbf{Training dynamics.}
Learning proceeds in four distinct phases: random seeding (eps 1--500),
rapid policy crystallisation (501--1000), convergence (1001--3000),
and plateau (3001--10000), with the zero-episode fraction dropping
from 20.8\% to 0.2\% in Phase~II alone.

\textbf{Q-value semantics.}
Decoding the highest and lowest Q-table entries reveals coherent
navigation logic: food-aligned, danger-free states receive Q-values
of 21--24; triply-trapped states receive Q-values of $-10$.
This confirms that the tabular agent has not merely memorised
reward patterns but has learned generalisable collision-avoidance
and food-seeking knowledge within the reachable subspace.

\textbf{Exploration scheduling dominates algorithm selection.}
The spread between exploration-scheduling extremes (14.79 for slow
decay vs 12.79 for low initial $\varepsilon$) is 2.0 points,
dwarfing the spread between algorithm families (13.73 Q-learning
vs 13.33 SARSA, gap 0.40 points).

\textbf{Reward shaping: a cautionary tale.}
Simple shaping accelerates early training (MA$+$15.6\% at ep500)
but reduces final greedy policy quality ($-5.4\%$), a clear case
of reward hacking.
Potential-based shaping avoids this by preserving policy invariance,
yielding a small but consistent $+0.5\%$ improvement with lower
score variance.

\textbf{Structural ceiling.}
By the end of training, all variants saturate in a narrow range of
256--259 reachable states, and the training MA score never exceeds 13
for the main run due to the one-step danger horizon of the 11-bit
encoding.
Overcoming this ceiling requires either richer state representations
or deep function approximation, motivating the extension to DQN
and POMDP-aware architectures for future work.

%% -------------------------------------------------------------
\bibliographystyle{IEEEtran}
\begin{thebibliography}{99}

\bibitem{watkins1992q}
C.~J.~C.~H.\ Watkins and P.\ Dayan,
``Q-learning,''
\emph{Machine Learning}, vol.~8, no.~3--4, pp.~279--292, 1992.

\bibitem{sutton2018reinforcement}
R.~S.\ Sutton and A.~G.\ Barto,
\emph{Reinforcement Learning: An Introduction}, 2nd~ed.
Cambridge, MA, USA: MIT Press, 2018.

\bibitem{mnih2015human}
V.~Mnih \emph{et al.},
``Human-level control through deep reinforcement learning,''
\emph{Nature}, vol.~518, no.~7540, pp.~529--533, 2015.

\bibitem{silver2016mastering}
D.~Silver \emph{et al.},
``Mastering the game of Go with deep neural networks and tree search,''
\emph{Nature}, vol.~529, no.~7587, pp.~484--489, 2016.

\bibitem{vinyals2019grandmaster}
O.~Vinyals \emph{et al.},
``Grandmaster level in StarCraft~II using multi-agent reinforcement
learning,''
\emph{Nature}, vol.~575, no.~7782, pp.~350--354, 2019.

\bibitem{puterman1994markov}
M.~L.\ Puterman,
\emph{Markov Decision Processes: Discrete Stochastic Dynamic Programming}.
Hoboken, NJ, USA: Wiley, 1994.

\bibitem{kaelbling1998planning}
L.~P.\ Kaelbling, M.~L.\ Littman, and A.~R.\ Cassandra,
``Planning and acting in partially observable stochastic domains,''
\emph{Artif.\ Intell.}, vol.~101, no.~1--2, pp.~99--134, 1998.

\bibitem{rummery1994line}
G.~A.\ Rummery and M.\ Niranjan,
``On-line Q-learning using connectionist systems,''
Tech.\ Rep.\ CUED/F-INFENG/TR 166, Univ.\ of Cambridge, 1994.

\bibitem{hasselt2010double}
H.~van Hasselt,
``Double Q-learning,''
in \emph{Adv.\ Neural Inf.\ Process.\ Syst.\ (NeurIPS)},
vol.~23, pp.~2613--2621, 2010.

\bibitem{ng1999policy}
A.~Y.\ Ng, D.\ Harada, and S.\ Russell,
``Policy invariance under reward transformations: Theory and application
to reward shaping,''
in \emph{Proc.\ ICML}, pp.~278--287, 1999.

\bibitem{auer2002finite}
P.~Auer, N.\ Cesa-Bianchi, and P.\ Fischer,
``Finite-time analysis of the multiarmed bandit problem,''
\emph{Machine Learning}, vol.~47, no.~2--3, pp.~235--256, 2002.

\bibitem{schaul2015prioritized}
T.~Schaul, J.\ Quan, I.\ Antonoglou, and D.\ Silver,
``Prioritized experience replay,''
in \emph{Proc.\ ICLR}, 2016.

\bibitem{barto2003recent}
A.~G.\ Barto and S.\ Mahadevan,
``Recent advances in hierarchical reinforcement learning,''
\emph{Discrete Event Dyn.\ Syst.}, vol.~13, no.~4, pp.~341--379, 2003.

\bibitem{sutton1999between}
R.~S.\ Sutton, D.\ Precup, and S.\ Singh,
``Between MDPs and semi-MDPs: A framework for temporal abstraction
in reinforcement learning,''
\emph{Artif.\ Intell.}, vol.~112, no.~1--2, pp.~181--211, 1999.

\bibitem{taylor2009transfer}
M.~E.\ Taylor and P.\ Stone,
``Transfer learning for reinforcement learning domains: A survey,''
\emph{J.\ Mach.\ Learn.\ Res.}, vol.~10, pp.~1633--1685, 2009.

\bibitem{finn2017model}
C.~Finn, P.\ Abbeel, and S.\ Levine,
``Model-agnostic meta-learning for fast adaptation of deep networks,''
in \emph{Proc.\ ICML}, pp.~1126--1135, 2017.

\bibitem{bellemare2016unifying}
M.~G.\ Bellemare \emph{et al.},
``Unifying count-based exploration and intrinsic motivation,''
in \emph{Adv.\ Neural Inf.\ Process.\ Syst.\ (NeurIPS)},
vol.~29, pp.~1471--1479, 2016.

\bibitem{pathak2017curiosity}
D.~Pathak, P.\ Agrawal, A.~A.\ Efros, and T.\ Darrell,
``Curiosity-driven exploration by self-supervised prediction,''
in \emph{Proc.\ ICML}, pp.~2778--2787, 2017.

\bibitem{kapturowski2018recurrent}
S.~Kapturowski, G.\ Ostrovski, J.\ Quan, R.\ Munos,
and W.\ Dabney,
``Recurrent experience replay in distributed reinforcement learning,''
in \emph{Proc.\ ICLR}, 2019.

\bibitem{ring1994continual}
M.~B.\ Ring,
``Continual learning in reinforcement learning systems,''
Ph.D.\ dissertation, Univ.\ of Texas at Austin, 1994.

\end{thebibliography}

\end{document}
